\subsection{Value-at-Risk and Expected-Shortfall Back-Test (Pipeline~A)}
\label{sec:var_backtest}
\label{sec:utility}

Still under Pipeline~A (Figure~\ref{fig:pipeline_schematic}): VaR and Expected-Shortfall are evaluated on the per-ticker SPY simulations, no cross-asset coupling.
For a financial-risk consumer the canonical diagnostics are one-tail Value-at-Risk (VaR) and Expected-Shortfall (ES) at daily probability levels of $1\%$ and $5\%$.
A synthetic-data generator is well calibrated if the envelope of per-path simulated VaR/ES brackets the observed historical value.
We simulate $1{,}000$ independent paths from each of the six most relevant generators (Bootstrap, GARCH(1,1), CHMM-N, CHMM-t, CHMM-L, CHMM-GED) and report median and $[5\%, 95\%]$ envelope across paths in Table~\ref{tab:var_es}.

\begin{table}[t]
\centering
\caption{\textbf{VaR and Expected-Shortfall envelope for SPY} ($N_{\text{paths}} = 1{,}000$). Each entry is the median and $[5\%, 95\%]$ envelope across paths of the left-tail VaR / ES at $\alpha \in \{0.01, 0.05\}$; observed historical values are in the first row. All four CHMM variants bracket the observed historical VaR and ES at both levels on both windows. Note on the repeated $\text{LR}_{\text{uc}} = 3.83$ values across CHMM rows in the body discussion: the unconditional Kupiec statistic at $T = 572$ trading days and $p = 0.05$ takes integer-breach-count values (e.g., $k = 19$ breaches gives $\text{LR}_{\text{uc}} = 3.828$, just below the asymptotic $\chi^2_1$ critical value of $3.841$); the apparent agreement across rows reflects the coarse integer-breach grid at this $T_{\text{OoS}}$, not a numerical coincidence in the underlying simulated samples.}
\label{tab:var_es}
\small
\resizebox{\textwidth}{!}{%
\begin{tabular}{l l l l l}
\toprule
          & IS VaR$_{0.01}$              & IS ES$_{0.01}$               & IS VaR$_{0.05}$              & IS ES$_{0.05}$               \\
Model     & median $[5\%, 95\%]$         & median $[5\%, 95\%]$         & median $[5\%, 95\%]$         & median $[5\%, 95\%]$         \\
\midrule
\textit{Observed}  & $-6.38$                      & $-9.00$                      & $-3.43$                      & $-5.50$                      \\
Bootstrap & $-6.38\ [-7.04, -5.99]$      & $-8.86\ [-10.37,\ -7.78]$    & $-3.43\ [-3.68, -3.20]$      & $-5.46\ [-5.99, -5.05]$      \\
GARCH     & $-5.80\ [-8.62, -4.52]$      & $-7.86\ [-13.54,\ -5.77]$    & $-3.19\ [-3.92, -2.71]$      & $-4.89\ [-7.11, -3.91]$      \\
CHMM-N    & $-6.64\ [-8.01, -5.47]$      & $-8.80\ [-10.16,\ -7.34]$    & $-3.45\ [-4.05, -2.95]$      & $-5.45\ [-6.33, -4.61]$      \\
CHMM-t    & $-6.58\ [-7.51, -5.83]$      & $-9.06\ [-10.70,\ -7.79]$    & $-3.40\ [-3.90, -2.90]$      & $-5.51\ [-6.18, -4.85]$      \\
CHMM-L    & $-6.89\ [-7.85, -6.01]$      & $-9.15\ [-10.42,\ -7.97]$    & $-3.30\ [-3.77, -2.88]$      & $-5.53\ [-6.17, -4.87]$      \\
CHMM-GED  & $-6.85\ [-7.70, -6.06]$      & $-8.79\ [-9.83,\ -7.82]$     & $-3.43\ [-3.98, -2.93]$      & $-5.55\ [-6.23, -4.91]$      \\
\midrule
          & OoS VaR$_{0.01}$             & OoS ES$_{0.01}$              & OoS VaR$_{0.05}$             & OoS ES$_{0.05}$              \\
\textit{Observed}  & $-5.51$                      & $-7.94$                      & $-2.77$                      & $-4.67$                      \\
Bootstrap & $-6.33\ [-7.60, -5.28]$      & $-8.43\ [-11.76,\ -6.59]$    & $-3.39\ [-3.91, -2.88]$      & $-5.34\ [-6.38, -4.50]$      \\
GARCH     & $-5.35\ [-9.98, -3.61]$      & $-6.72\ [-13.51,\ -4.38]$    & $-3.14\ [-4.94, -2.36]$      & $-4.59\ [-7.92, -3.23]$      \\
CHMM-N    & $-6.30\ [-9.02, -4.30]$      & $-8.44\ [-11.36,\ -5.42]$    & $-3.39\ [-4.74, -2.47]$      & $-5.27\ [-7.23, -3.74]$      \\
CHMM-t    & $-6.32\ [-8.13, -4.80]$      & $-8.60\ [-12.20,\ -6.29]$    & $-3.29\ [-4.27, -2.45]$      & $-5.34\ [-6.72, -4.08]$      \\
CHMM-L    & $-6.69\ [-8.59, -4.91]$      & $-8.88\ [-11.48,\ -6.73]$    & $-3.27\ [-4.30, -2.51]$      & $-5.47\ [-6.82, -4.20]$      \\
CHMM-GED  & $-6.60\ [-8.30, -4.82]$      & $-8.51\ [-10.68,\ -6.47]$    & $-3.32\ [-4.38, -2.46]$      & $-5.41\ [-6.66, -4.06]$      \\
\bottomrule
\end{tabular}%
}
\end{table}

All four CHMM variants bracket the observed historical VaR and ES at both $1\%$ and $5\%$ levels within their $5\%$--$95\%$ envelopes on both IS and OoS windows; the IS CHMM-N median VaR$_{0.01}$ at $-6.64$ (observed $-6.38$) and ES$_{0.01}$ at $-8.80$ (observed $-9.00$) are the closest point estimates across the six generators.
The CHMM envelopes are roughly $1.7\times$ tighter than GARCH's and pass the unconditional Kupiec test cleanly at both levels (CHMM-N $\text{LR}_{\text{uc}} = 3.83$ at $5\%$ against the $\chi^2_1 = 3.84$ critical value; CHMM-t at $3.83$; CHMM-L at $3.03$; CHMM-GED at $3.83$).
CHMM-t's IS kurtosis overshoot does \emph{not} translate into a systematically more conservative VaR/ES point estimate; the IS median VaR$_{0.01}$ of $-6.58$ is almost identical to CHMM-N's $-6.64$ and the envelope width is comparable. CHMM-L's IS median VaR$_{0.01}$ of $-6.89$ is mildly more conservative than CHMM-N's, consistent with its closer-to-observed kurtosis. CHMM-GED's IS median VaR$_{0.01}$ of $-6.85$ sits between CHMM-N and CHMM-L on the conservativeness axis, consistent with its adaptive shape mechanism.

\paragraph{Conditional coverage (Christoffersen 1998).}
The Kupiec statistic only checks the unconditional breach rate. We complement it with the Christoffersen \cite{christoffersen1998evaluating} independence ($\text{LR}_{\text{ind}}$, $\sim\chi^2_1$) and joint conditional-coverage ($\text{LR}_{\text{cc}} = \text{LR}_{\text{uc}} + \text{LR}_{\text{ind}}$, $\sim\chi^2_2$) tests. We report two variants: the unconditional pooled-archive variant matched to the rest of the body panel, and a regime-conditional variant that propagates the latent-state forecast $\Prob(s_{t+1}\,|\,\mathcal F_t)$ through the CHMM emission mixture under IS-fixed parameters. Every unconditional generator in the panel (i.i.d.\ bootstrap, GARCH(1,1), and all four CHMM variants) rejects independence on the OoS window at $\alpha = 0.05$ for both VaR levels (CHMM-N at $5\%$: $\text{LR}_{\text{ind}} = 5.26$, $p = 0.022$; $\text{LR}_{\text{cc}} = 9.08$, $p = 0.011$; full unconditional panel in Appendix~\ref{sec:christoffersen_supp}, Table~\ref{tab:christoffersen_var}). The independence rejection is driven by volatility-clustering in $R_{\text{oos}}$ that any unconditional generator with a constant-across-$t$ VaR threshold is structurally unable to track. The regime-conditional variant addresses this directly: at each OoS day $t$, we run the forward filter through $\mathcal F_{t-1} = R_{\text{IS}} \cup R_{\text{oos}}[1{:}t-1]$ to get the one-step-ahead state forecast $\Prob(s_t = k \mid \mathcal F_{t-1}) = \sum_i \Prob(s_{t-1} = i \mid \mathcal F_{t-1}) \cdot T_{ik}$ and define $\widehat{\text{VaR}}_t(\alpha) = F_t^{-1}(\alpha)$ where $F_t$ is the $K$-component Gaussian-mixture predictive CDF $\sum_k \Prob(s_t = k \mid \mathcal F_{t-1}) \cdot \Phi(\,\cdot\,;\, \mu_k, \sigma_k)$. Table~\ref{tab:cond_var} reports the resulting back-test for CHMM-N at $K^\star = 3$ and at $K = 18$ on the OoS window.

\begin{table}[h]
\centering
\small
\caption{\textbf{Regime-conditional VaR back-test for CHMM-N on the OoS window} ($T_{\text{oos}} = 572$ days, seed $= 20260420$). At each $t$, conditional VaR is the $\alpha$-quantile of the one-step-ahead Gaussian-mixture predictive density under IS-fixed parameters. Critical values: $\chi^2_1(0.05) = 3.841$, $\chi^2_2(0.05) = 5.991$. Every $(K, \alpha)$ row passes Kupiec, Christoffersen-ind, and Christoffersen-cc cleanly at $\alpha = 0.05$, in contrast to the unconditional pooled-archive variant which rejects independence at every CHMM, GARCH, and bootstrap row on OoS (Table~\ref{tab:christoffersen_var}).}
\label{tab:cond_var}
\begin{tabular}{c c c c c c c c c c}
\toprule
$K$ & $\alpha$ & breaches & br rate & med VaR & $\text{LR}_{\text{uc}}$ & $p_{\text{uc}}$ & $\text{LR}_{\text{ind}}$ & $\text{LR}_{\text{cc}}$ & $p_{\text{cc}}$ \\
\midrule
$3$  & $0.01$ & $9$  & $1.57\%$ & $-4.56$ & $1.62$ & $0.20$ & $2.36$ & $3.98$ & $0.14$ \\
$3$  & $0.05$ & $35$ & $6.12\%$ & $-2.87$ & $1.41$ & $0.24$ & $0.01$ & $1.42$ & $0.49$ \\
$18$ & $0.01$ & $9$  & $1.57\%$ & $-5.20$ & $1.62$ & $0.20$ & $2.36$ & $3.98$ & $0.14$ \\
$18$ & $0.05$ & $26$ & $4.55\%$ & $-3.02$ & $0.26$ & $0.61$ & $0.52$ & $0.78$ & $0.68$ \\
\bottomrule
\end{tabular}
\end{table}

The conditional-VaR construction passes all three tests at every $(K, \alpha)$ on OoS, with breach rates of $1.57\%$ and $4.55$--$6.12\%$ against nominal $1\%$ and $5\%$. The independence-test improvement from the unconditional variant ($\text{LR}_{\text{ind}} = 5.26$ at $K = 18$, $\alpha = 0.05$) to the conditional variant ($\text{LR}_{\text{ind}} = 0.52$) is the regime-switching value proposition that the Kupiec-only headline did not exercise: the latent-state forecast tightens the threshold in clustered-volatility windows so that breaches do not pile up at sub-runs of high volatility. The same construction does not require retuning $K$ on a stress-test slice nor a different parameter family for VaR consumers; it reuses the IS-fixed CHMM-N transitions and emissions and runs the forward filter on the OoS series under those frozen parameters. The full unconditional Christoffersen panel across the six panel rows (Bootstrap, GARCH, CHMM-N/-t/-L/-GED) is in Appendix~\ref{sec:christoffersen_supp}.
